\paragraph{Leitura.} A penalidade de TS2-D aparece em todas as classes que leem intervalos (1,60--1,80$\times$), porque o bucket diário multiplica as chamadas CQL (uma por dia da janela de 7 dias). Q1 (ponto) e Q8 (último valor) não mudam. TS0 paga uma partição por sensor: Q1 sai 1,96$\times$ mais lento e Q8 1,72$\times$. A referência original só perde no Q8 (24,9$\times$, varredura completa); com o Q8 otimizado ela é indistinguível de TS2-U em todas as classes. \textbf{Observação metodológica:} no par A/A REF-Q8opt $\equiv$ TS2-U as razões por classe variam de 0,91 a 1,05 (por exemplo, Q4), embora o trabalho seja idêntico. Por isso diferenças de até $\approx$17\% numa classe isolada não são interpretáveis, e a decisão é tomada no mix, cujo piso é de 2,2\%.
